\documentclass[11pt,a4paper]{article}

\usepackage[utf8]{inputenc}
\usepackage[T1]{fontenc}
\usepackage[french]{babel}
\usepackage[margin=2.5cm]{geometry}
\usepackage{graphicx}
\usepackage{float}
\usepackage{booktabs}
\usepackage{amsmath}
\usepackage{listings}
\usepackage{hyperref}

% --- listings (code C avec accents) ---
\lstset{
  language=C,
  basicstyle=\ttfamily\footnotesize,
  breaklines=true,
  frame=single,
  numbers=left,
  numberstyle=\tiny,
  inputencoding=utf8,
  extendedchars=true,
  literate={é}{{\'e}}1 {è}{{\`e}}1 {à}{{\`a}}1 {É}{{\'E}}1 {ê}{{\^e}}1 {ç}{{\c c}}1
}

% --- emplacement de capture d'ecran ---
% Pour inserer une vraie capture, remplacer l'appel \capture{...}{...}
% par : \begin{figure}[H]\centering
%       \includegraphics[width=\textwidth]{captures/nom.png}
%       \caption{...}\end{figure}
\newcommand{\capture}[2]{%
  \begin{figure}[H]
    \centering
    \fbox{\parbox[c][#1][c]{0.92\textwidth}{\centering\textit{Emplacement de la capture d'écran}}}
    \caption{#2}
  \end{figure}}

\title{INFO932 -- TP1\\Programmation parallèle avec OpenMP}
\author{Prénom NOM}
\date{\today}

\begin{document}
\maketitle

\section{Introduction}

L'objectif de ce TP est d'écrire un programme séquentiel en C qui manipule un grand tableau d'entiers, puis de le paralléliser avec OpenMP et de comparer les performances.

Le programme réalise les six opérations suivantes :
\begin{enumerate}
  \item allocation d'un grand tableau et remplissage avec des entiers aléatoires ;
  \item recherche du maximum ;
  \item recherche du minimum ;
  \item calcul de la somme des éléments ;
  \item comptage des éléments compris entre deux bornes (le minimum et le maximum) ;
  \item tri du tableau par tri sélection.
\end{enumerate}

Nous mesurons le temps de chaque étape avec un nombre variable de threads afin de comprendre ce qui se parallélise bien, ce qui se parallélise mal, et pourquoi augmenter le nombre de c\oe urs n'accélère pas toujours le programme.

\section{Environnement d'exécution}

Les mesures ont été faites sur la machine décrite ci-dessous. Le nombre de c\oe urs physiques, le nombre de threads matériels (hyper-threading) et le système d'exploitation sont déterminants pour interpréter les résultats.

\capture{5cm}{Caractéristiques de la machine (\texttt{lscpu}, \texttt{nproc}, \texttt{gcc --version}).}

\begin{table}[H]
\centering
\begin{tabular}{ll}
\toprule
Processeur & \ldots \\
C\oe urs physiques & \ldots \\
Threads matériels (logiques) & \ldots \\
Mémoire vive & \ldots \\
Système d'exploitation & \ldots \\
Compilateur & gcc \ldots{} (options \texttt{-O2}, \texttt{-fopenmp}) \\
\bottomrule
\end{tabular}
\caption{Environnement de test.}
\end{table}

\section{Programme séquentiel}

\subsection{Choix de conception}

Le programme (\texttt{sequentiel.c}) reçoit la taille $N$ du tableau en argument et l'alloue avec \texttt{malloc}. Chaque étape est implémentée dans une fonction séparée : \texttt{remplirTableau}, \texttt{maximum}, \texttt{minimum}, \texttt{somme}, \texttt{between} et \texttt{triSelection}.

\begin{itemize}
  \item La somme est stockée dans un \texttt{long long} pour éviter tout débordement sur les grandes tailles.
  \item La fonction \texttt{between} compte les éléments \emph{strictement} compris entre le minimum et le maximum du tableau. Avec des bornes inclusives, le résultat vaudrait toujours $N$, puisque tout élément est compris entre le minimum et le maximum.
  \item Chaque étape est chronométrée séparément avec \texttt{clock\_gettime(CLOCK\_MONOTONIC)}.
  \item Le tri sélection a une complexité en $O(N^2)$. Il n'est exécuté que pour $N \leq 200\,000$ ; au-delà, le temps d'exécution devient prohibitif.
\end{itemize}

\subsection{Complexité}

Les étapes 1 à 5 sont en $O(N)$ : un seul parcours du tableau. Le tri sélection est en $O(N^2)$ : pour $N = 100\,000$, il effectue environ $5 \times 10^9$ comparaisons. Il domine donc largement le temps total.

\section{Parallélisation avec OpenMP}

Le programme parallèle (\texttt{parallele.c}) conserve la même structure. Les modifications sont résumées dans le tableau suivant.

\begin{table}[H]
\centering
\small
\begin{tabular}{p{3.2cm}p{4.3cm}p{6.5cm}}
\toprule
Fonction & Directive OpenMP & Justification \\
\midrule
\texttt{remplirTableau} & \texttt{parallel} + \texttt{for} & \texttt{rand()} n'est pas thread-safe ; chaque thread utilise \texttt{rand\_r()} avec sa propre graine. \\
\texttt{maximum} & \texttt{parallel for reduction(max)} & Chaque thread calcule un maximum local, combiné à la fin. \\
\texttt{minimum} & \texttt{parallel for reduction(min)} & Même principe. \\
\texttt{somme} & \texttt{parallel for reduction(+)} & Sommes partielles additionnées à la fin. \\
\texttt{between} & \texttt{parallel for reduction(+)} & Compteurs partiels additionnés à la fin. \\
\texttt{triSelection} & \texttt{parallel}, \texttt{for nowait}, \texttt{critical}, \texttt{barrier}, \texttt{single} & Seule la recherche du minimum (boucle interne) est parallélisée. \\
\bottomrule
\end{tabular}
\caption{Parallélisation de chaque fonction.}
\end{table}

\subsection{Les réductions}

Sans la clause \texttt{reduction}, plusieurs threads modifieraient en même temps la même variable (condition de course) et le résultat serait faux. Avec \texttt{reduction}, chaque thread travaille sur une copie privée, et OpenMP combine ces copies à la fin de la boucle. Ces quatre fonctions sont dites \emph{embarrassingly parallel} : les itérations sont indépendantes.

\subsection{Le tri sélection}

Le tri sélection est difficile à paralléliser : l'itération $i$ de la boucle externe dépend de l'état du tableau après l'itération $i-1$. La boucle externe reste donc séquentielle. Nous parallélisons la boucle interne, qui cherche le minimum de la partie non triée :
\begin{enumerate}
  \item un thread initialise le minimum partagé (\texttt{single}) ;
  \item chaque thread cherche le minimum dans sa tranche du tableau (\texttt{for nowait}) ;
  \item les minimums locaux sont fusionnés dans une section \texttt{critical} ;
  \item une barrière (\texttt{barrier}) garantit que tous les threads ont fini de fusionner ;
  \item un thread effectue l'échange (\texttt{single}).
\end{enumerate}

La région parallèle est créée une seule fois, en dehors de la boucle externe, pour éviter de recréer les threads à chaque itération. En contrepartie, chaque itération impose deux synchronisations (la barrière et la fin du \texttt{single}) : pour $N = 100\,000$, cela représente environ $200\,000$ synchronisations. Ce coût est la clé de l'analyse du tri.

\section{Protocole expérimental}

\subsection{Compilation}

\begin{verbatim}
gcc -O2 sequentiel.c -o runnable/sequentiel
gcc -fopenmp -O2 parallele.c -o runnable/parallele
\end{verbatim}

\subsection{Exécutions}

Deux tailles sont testées : $N = 100\,000$ (tri inclus, étapes $O(N)$ très courtes) et $N = 50\,000\,000$ (étapes $O(N)$ mesurables, tri ignoré). Le nombre de threads $p$ est fixé avec \texttt{OMP\_NUM\_THREADS}.

\begin{verbatim}
./runnable/sequentiel 100000
./runnable/sequentiel 50000000
for t in 1 2 4 8; do OMP_NUM_THREADS=$t ./runnable/parallele 100000; done
for t in 1 2 4 8; do OMP_NUM_THREADS=$t ./runnable/parallele 50000000; done
\end{verbatim}

\subsection{Métriques}

Pour une étape donnée, avec $T_s$ le temps séquentiel et $T_p$ le temps avec $p$ threads :
\[
S(p) = \frac{T_s}{T_p} \qquad\text{(accélération)}, \qquad
E(p) = \frac{S(p)}{p} \qquad\text{(efficacité)}.
\]
Une accélération idéale vaut $S(p) = p$, soit une efficacité de $1$. Une accélération inférieure à $1$ signifie que la version parallèle est \emph{plus lente} que la version séquentielle.

\section{Résultats}

\subsection{Version séquentielle}

\capture{6cm}{Exécution séquentielle pour $N = 100\,000$.}

\capture{6cm}{Exécution séquentielle pour $N = 50\,000\,000$.}

\subsection{Version parallèle, $N = 100\,000$}

\capture{9cm}{Exécutions parallèles pour $N = 100\,000$ avec 1, 2, 4 et 8 threads.}

\begin{table}[H]
\centering
\small
\begin{tabular}{lrrrrr}
\toprule
Étape (s) & Séquentiel & 1 thread & 2 threads & 4 threads & 8 threads \\
\midrule
Remplissage & \ldots & \ldots & \ldots & \ldots & \ldots \\
Maximum & \ldots & \ldots & \ldots & \ldots & \ldots \\
Minimum & \ldots & \ldots & \ldots & \ldots & \ldots \\
Somme & \ldots & \ldots & \ldots & \ldots & \ldots \\
Between & \ldots & \ldots & \ldots & \ldots & \ldots \\
Tri sélection & \ldots & \ldots & \ldots & \ldots & \ldots \\
\midrule
Total & \ldots & \ldots & \ldots & \ldots & \ldots \\
\bottomrule
\end{tabular}
\caption{Temps d'exécution par étape pour $N = 100\,000$.}
\end{table}

\begin{table}[H]
\centering
\begin{tabular}{lrrrr}
\toprule
Tri sélection & 1 thread & 2 threads & 4 threads & 8 threads \\
\midrule
Accélération $S(p)$ & \ldots & \ldots & \ldots & \ldots \\
Efficacité $E(p)$ & \ldots & \ldots & \ldots & \ldots \\
\bottomrule
\end{tabular}
\caption{Accélération et efficacité du tri sélection ($N = 100\,000$).}
\end{table}

\subsection{Version parallèle, $N = 50\,000\,000$}

\capture{9cm}{Exécutions parallèles pour $N = 50\,000\,000$ avec 1, 2, 4 et 8 threads.}

\begin{table}[H]
\centering
\small
\begin{tabular}{lrrrrr}
\toprule
Étape (s) & Séquentiel & 1 thread & 2 threads & 4 threads & 8 threads \\
\midrule
Remplissage & \ldots & \ldots & \ldots & \ldots & \ldots \\
Maximum & \ldots & \ldots & \ldots & \ldots & \ldots \\
Minimum & \ldots & \ldots & \ldots & \ldots & \ldots \\
Somme & \ldots & \ldots & \ldots & \ldots & \ldots \\
Between & \ldots & \ldots & \ldots & \ldots & \ldots \\
\bottomrule
\end{tabular}
\caption{Temps d'exécution par étape pour $N = 50\,000\,000$.}
\end{table}

\begin{table}[H]
\centering
\begin{tabular}{lrrrr}
\toprule
Accélération $S(p)$ & 1 thread & 2 threads & 4 threads & 8 threads \\
\midrule
Remplissage & \ldots & \ldots & \ldots & \ldots \\
Maximum & \ldots & \ldots & \ldots & \ldots \\
Minimum & \ldots & \ldots & \ldots & \ldots \\
Somme & \ldots & \ldots & \ldots & \ldots \\
Between & \ldots & \ldots & \ldots & \ldots \\
\bottomrule
\end{tabular}
\caption{Accélération par étape pour $N = 50\,000\,000$.}
\end{table}

\section{Analyse}

\subsection{Un processeur contre plusieurs}

Avec un seul thread, la version OpenMP est légèrement plus lente que la version séquentielle : la création de la région parallèle et la gestion des réductions ont un coût, même sans parallélisme réel. Avec plusieurs threads, le travail est réparti, mais l'accélération dépend de la taille du problème.

\paragraph{Petit tableau ($N = 100\,000$).} Les étapes en $O(N)$ ne durent que quelques dizaines de microsecondes. Le coût de création et de synchronisation des threads (de l'ordre de la dizaine de microsecondes) est du même ordre que le calcul lui-même. Il est donc normal que la version parallèle soit aussi rapide, voire plus lente, que la version séquentielle pour ces étapes.

\paragraph{Grand tableau ($N = 50\,000\,000$).} Le calcul domine le coût de gestion des threads et l'accélération devient visible. Elle reste toutefois inférieure à l'idéal $S(p) = p$, car ces opérations sont limitées par la \emph{bande passante mémoire} : chaque élément est lu une seule fois pour très peu de calcul, et la mémoire ne peut pas alimenter indéfiniment un nombre croissant de c\oe urs.

\subsection{Pourquoi 4 c\oe urs peuvent être plus rapides que 8}

Il est tout à fait possible, et nos mesures le montrent, que l'exécution avec 4 threads soit plus rapide qu'avec 8. Plusieurs raisons se cumulent.

\begin{enumerate}
  \item \textbf{Cœurs physiques et hyper-threading.} Si la machine possède 4 c\oe urs physiques et 8 threads logiques, deux threads se partagent les unités de calcul et le cache d'un même c\oe ur. Les 8 threads n'apportent alors pas deux fois plus de puissance de calcul que 4.
  \item \textbf{Surcoût de synchronisation.} Le tri sélection impose deux barrières par itération. Le coût d'une barrière augmente avec le nombre de threads, car il faut attendre le plus lent. Avec 8 threads, la durée de la synchronisation peut dépasser le gain obtenu en divisant la tranche de recherche.
  \item \textbf{Travail par thread trop faible.} Dans le tri sélection, la tranche à parcourir diminue à chaque itération ($N - i$ éléments). Vers la fin du tri, elle ne contient plus que quelques éléments par thread : le travail utile est négligeable devant la synchronisation.
  \item \textbf{Section critique.} Les $p$ threads se retrouvent dans la section \texttt{critical} à chaque itération et passent l'un après l'autre. Son coût croît donc avec $p$.
  \item \textbf{Sursouscription.} Si le nombre de threads dépasse le nombre de c\oe urs disponibles, le système d'exploitation doit les ordonnancer à tour de rôle. Un thread en attente à une barrière peut alors bloquer des threads qui avaient du travail à faire.
  \item \textbf{Bande passante mémoire saturée.} Pour les parcours simples, 4 threads peuvent déjà saturer la mémoire : en ajouter ne fait qu'augmenter la contention.
\end{enumerate}

Cet effet illustre la loi d'Amdahl : si $f$ est la fraction parallélisable du programme,
\[
S(p) = \frac{1}{(1-f) + f/p}.
\]
L'accélération est bornée par la partie séquentielle, et le modèle ne tient pas compte du surcoût de synchronisation, qui \emph{croît} avec $p$. L'efficacité $E(p)$ diminue donc quand $p$ augmente, et au-delà d'un certain nombre de threads le temps total peut même remonter.

\subsection{Analyse des mesures}

\emph{[À compléter avec les résultats obtenus : valeur des accélérations, nombre de threads optimal pour le tri et pour les étapes $O(N)$, comparaison 4 threads / 8 threads.]}

\section{Conclusion}

La parallélisation avec OpenMP est très simple à mettre en \oe uvre pour les boucles indépendantes (remplissage, maximum, minimum, somme, comptage) grâce aux directives \texttt{parallel for} et \texttt{reduction}, mais le gain dépend fortement de la taille du problème : sur un petit tableau, le coût de gestion des threads annule le bénéfice. Le tri sélection, qui contient une dépendance entre itérations, ne se parallélise que partiellement et son accélération est limitée par les synchronisations. Enfin, utiliser plus de threads n'est pas toujours plus rapide : le nombre de threads optimal dépend du nombre de c\oe urs physiques, de la taille des données et du ratio entre calcul et synchronisation.

\appendix

\section{Code source}

\subsection{Programme séquentiel}
\lstinputlisting{sequentiel.c}

\subsection{Programme parallèle}
\lstinputlisting{parallele.c}

\end{document}
